%% kapa: exam=1 round=21 year=2010 period=1 form=A subject=civil_law subject_ko=민법 seq=17 number=17
%% source: Q-Net 감정평가사 기출문제 (cst00309) artlSeq=5024181 file=extracted/1008/5024181/2010년 제21회 감정평가사 1차시험 기출문제/1교시.hwp
%% conversion: kordoc 4.24.51 -> tools/gen1.py (auto); review=proofread-sonnet; flags=none
\begin{question}{17}
부동산의 취득시효에 관한 설명으로 옳지 않은 것은? (다툼이 있으면 판례에 의함)
\begin{choices}
  \item 국유의 일반재산(잡종재산)에 대하여 취득시효가 완성된 후에 그 일반재산이 행정재산으로 되었다면 시효완성을 이유로 소유권이전등기를 청구할 수 없다.
  \item 점유취득시효 완성 후 제3자 명의의 소유권이전등기가 마쳐진 경우, 점유자는 그 소유권 변동시를 새로운 기산점으로 하여 2차의 취득시효의 완성을 주장할 수 있다.
  \item 시효기간 진행 중 토지의 점유자가 그 토지에 대하여 매매를 원인으로 한 소유권이전등기 청구소송을 제기하였다가 패소 확정되더라도, 점유자의 점유가 타주점유로 전환되는 것이 아니다.
  \item 시효기간 진행 중 소유자가 바뀐 경우 시효를 주장하는 자는 임의로 시효기간의 기산점을 정할 수 없다.
  \item 시효기간 만료 당시의 점유자로부터 점유를 승계한 현 점유자가 전 점유자의 취득시효 완성의 효과를 주장하여 직접 자기에게 소유권이전등기를 청구할 수 있다.
\end{choices}
\end{question}
\answer{5}
